% =====================================================================
%  MLOps (UPC 2026-27 Q1) - Team Vera - Project report
%  Structure mirrors the official Atenea template (mlops_report_template.docx).
%  Grey italic text = template "Guideline" prompts to be replaced by the team.
%  Section 2.1.1 is filled with our finalised requirements-engineering content.
% =====================================================================
\documentclass[11pt]{article}
\usepackage[utf8]{inputenc}
\usepackage[T1]{fontenc}
\usepackage[margin=2.3cm]{geometry}
\usepackage{booktabs}
\usepackage{tabularx}
\usepackage{array}
\usepackage{enumitem}
\usepackage{amssymb}
\usepackage{graphicx}
\usepackage{xcolor}
\usepackage{hyperref}
\hypersetup{colorlinks=true,linkcolor=black,citecolor=black,urlcolor=blue}
\newcolumntype{Y}{>{\raggedright\arraybackslash}X}
\newcommand{\guideline}[1]{\par{\color{gray}\itshape #1}\par\smallskip}
\setlength{\tabcolsep}{5pt}
\renewcommand{\arraystretch}{1.2}

\begin{document}

% ============================ COVER PAGE ============================
\begin{titlepage}
\centering
\vspace*{1.5cm}
\includegraphics[width=0.55\textwidth]{images/upc-positiu-p3005.png}\par
\vspace{3cm}
{\large\scshape Machine Learning Systems in Production (MLOps)\par}
\vspace{0.3cm}
{\normalsize Master in Data Science\par}
\vspace{3cm}
{\Huge\bfseries Detecting Real vs.\ AI-Generated Images\par}
\vspace{0.5cm}
{\large\itshape Initial deliverable --- Milestones 1--3\par}
\vspace{3cm}
{\large\bfseries Team Vera\par}
\vspace{0.6cm}
{\normalsize Arman Bazarchi \quad Adri\'an Segura \quad Pablo Rodr\'iguez\par}
\vfill
\end{titlepage}

% ============================== INDEX ==============================
\tableofcontents
\newpage

% =========================== 1. INTRODUCTION =======================
\section{Introduction}

\subsection{Goal of the project}
Build, deploy, and operate an ML component that, given an input image, decides whether it is real or AI-generated, exposed as a monitored production service. Our main task is binary image classification (real vs \texttt{ai\_generated}), and optionally we might as well classify AI-generated images among the 5 generators (which generator model owns the image). Our focus is mostly on engineering around the model (reproducibility, QA, deployment, monitoring), not maximizing raw accuracy.

\paragraph{Success criteria.} The component is successful when it is served via a reproducible, containerised API that is tested and deployed automatically on every change (CI/CD), and meets the model targets on a leakage-free test set with shortcuts controlled, while all processes monitored for drift.

\subsection{Teammates' evaluation}
\guideline{Option 1: A complete team consensus on the following statement:}
``All team members agree that they had an equal contribution to this delivery and project: completing a fair share of the team's work with acceptable/high quality, keeping commitments, and completing assignments on time, helping teammates who are having difficulty when it is easy or important''.

\guideline{Option 2: There is no consensus on the above statement. Each team member evaluates (see scale 1 to 5 below) his/her peers' contributions in each delivery according to their contribution, based on CATME ratings.}

\smallskip
\noindent\textit{Figure 1. CATME ratings (from \url{https://info.catme.org/features/catme-five-dimensions/})}
\smallskip

\noindent\begin{tabularx}{\textwidth}{@{}c Y@{}}
\toprule
\textbf{Rating} & \textbf{Description of Rating} \\
\midrule
5 & Does more or higher-quality work than expected. Makes important contributions that improve the team's work. Helps teammates who are having difficulty completing their work. \\
4 & Demonstrates behaviors described immediately above and below. \\
3 & Completes a fair share of the team's work with acceptable quality. Keeps commitments and completes assignments on time. Helps teammates who are having difficulty when it is easy or important. \\
2 & Demonstrates behaviors described immediately above and below. \\
1 & Does not do a fair share of the team's work. Delivers sloppy or incomplete work. \\
\bottomrule
\end{tabularx}

\medskip
\guideline{Fill in the following table (received evaluations in rows):}
\begin{center}
\begin{tabular}{@{}lcccccc@{}}
\toprule
 & Name 1 & Name 2 & Name 3 & Name 4 & Name 5 & Average \\
\midrule
Name 1 & -- &    &    &    &    &  \\
Name 2 &    & -- &    &    &    &  \\
Name 3 &    &    & -- &    &    &  \\
Name 4 &    &    &    & -- &    &  \\
Name 5 &    &    &    &    & -- &  \\
\bottomrule
\end{tabular}
\end{center}

% =========================== 2. METHODOLOGY ========================

\newpage

\section{Methodology}

% ----- 2.1 Milestone 1 -----
\subsection{Milestone 1: Inception}

\subsubsection{Selection of problem and requirements engineering for ML}

\paragraph{Problem selection.} Modern text-to-image generators (Stable Diffusion, DALL-E, MidJourney, \dots) produce photorealistic images at scale, fuelling misinformation, fake reviews, and trust problems for platforms and media. Reliably telling a \textbf{real photograph} from an \textbf{AI-generated image} cannot be solved with hand-written rules, the cues are subtle, generator-specific, and evolving, which makes it a genuine machine-learning problem. We target real-vs-AI detection in particular because nowadays it is increasingly difficult to recognise whether a piece of media is real or AI-generated.

\paragraph{What makes it unique.} Unlike general image-classification tasks whose target classes are usually obvious to a human (weather, animals, traffic signs, \dots), we target a category that in the near future may be very hard to classify even for people --- telling a real image from an AI-generated one. Our component also covers \textbf{five recent} generators (2022--2024), and instead of being specific to one dataset, our component would be easily configured and used with different, bigger datasets.

\paragraph{Intended use.} As main use case of our ML component we can mention reliable platforms (such as news, social media, \dots) to distinguish real from fake (or generated) images to prevent misinformation and correctly flag likely AI-generated uploads via an API; another use case could be to quickly check the authenticity of an image with a confidence score, which is a well-documented and reproducible component for downstream developers to integrate in a larger system.

\paragraph{ML component specification (I/O contract).}
\begin{itemize}[nosep,leftmargin=1.4em]

\medskip

  \item \textbf{Input:} one image (JPEG/PNG/WEBP), bounded file size; preprocessed identically to training (fixed square resize + uniform re-encode).

\medskip

  \item \textbf{Output:} JSON \texttt{\{"label": "real" | "ai\_generated", "score": <float 0--1>\}}; optionally generator model as well.

\medskip

  \item \textbf{Operating context:} stateless REST service, containerised, with health and metrics endpoints.
\end{itemize}

\newpage

\paragraph{Requirements.} \textit{Priority:} \textbf{M} = Must, \textbf{S} = Should, \textbf{C} = Could.

\medskip\noindent\textbf{Functional (FR)}\par\smallskip
\noindent\begin{tabularx}{\textwidth}{@{}l Y c@{}}
\toprule
\textbf{ID} & \textbf{Requirement} & \textbf{Pri} \\
\midrule
FR-1 & Classify a single image as \texttt{real} / \texttt{ai\_generated} & M \\
FR-2 & Return a confidence score in [0,1] with each prediction & M \\
FR-3 & Expose prediction via a REST endpoint (\texttt{POST /predict}) returning JSON & M \\
FR-4 & Accept JPEG/PNG/WEBP within a size limit; reject invalid input with a clear error & M \\
FR-5 & Apply the \textbf{same} preprocessing at train and serve time (resize + re-encode) & M \\
FR-6 & Provide health/readiness and metrics endpoints & S \\
FR-7 & Optionally return the most-likely generator (attribution) & C \\
\bottomrule
\\

\end{tabularx}



\medskip\noindent\textbf{Data (DR)}\par\smallskip
\noindent\begin{tabularx}{\textwidth}{@{}l Y c@{}}
\toprule
\textbf{ID} & \textbf{Requirement} & \textbf{Pri} \\
\midrule
DR-1 & Use labelled real \& AI images from multiple recent generators & M \\
DR-2 & Data \& train/validation/test splits to be versioned and reproducible (DVC) & M \\
DR-3 & Handle class imbalance if observed in the data (via class weights or generator-balanced sampling) & M \\
DR-4 & Make sure no leakage such as image caption shared across train/val/test sets & M \\
DR-5 & Preprocessing to neutralise data to ensure the model does not rely on non-semantic \textbf{shortcuts} (such as aspect ratio, format) & M \\
DR-6 & Data usage complies with licence (CC BY 4.0, MS COCO) & M \\
DR-7 & Reserve a held-out generator (and data records) for drift evaluation & C \\
DR-8 & Automated data-quality checks (schema, class values, missingness, \dots) & S \\
\bottomrule

\\

\end{tabularx}

\medskip\noindent\textbf{Model / performance (MR)}\par\smallskip
\noindent\begin{tabularx}{\textwidth}{@{}l Y c@{}}
\toprule
\textbf{ID} & \textbf{Requirement} & \textbf{Pri} \\
\midrule
MR-1 & Balanced accuracy (target $\geq$ 0.85) & S \\
MR-2 & Macro-F1 and recall on \textbf{minority} class, usually real images (targets $\geq$ 0.85 / $\geq$ 0.80) & S \\
MR-3 & \textbf{PR-AUC} (target $\geq$ 0.90) & S \\
MR-4 & \textbf{Cross-generator} (leave-one-generator-out) balanced accuracy (target $\geq$ 0.70) & S \\
MR-5 & Decision threshold based on validation results (not fixed 0.5) & M \\
\bottomrule
\end{tabularx}

\smallskip
\noindent Accuracy alone is misleading if there is class imbalance in data, so success would be defined on balanced/minority-aware metrics and on generalisation to unseen generators.

\newpage


\medskip\noindent\textbf{Non-functional (NFR)}\par\smallskip


\noindent\begin{tabularx}{\textwidth}{@{}l Y c@{}}
\toprule
\textbf{ID} & \textbf{Requirement} & \textbf{Pri} \\
\midrule
NFR-1 & Runs without a GPU, on minimum 2 CPU cores and 8\,GB of RAM: retrieving the data (\texttt{dvc pull}) takes $\leq$ 10\,min, and processing it into a trained model (preprocess, split, embed, train) $\leq$ 10\,min with a peak memory $\leq$ 3\,GB & S \medskip \\
NFR-2 & With 10 concurrent clients, one container serves $\geq$ 5 requests/s for 5\,min with a p95 latency $\leq$ 2\,s and no failed requests; being stateless, it scales by adding replicas & S \medskip \\
NFR-3 & The code is modular and easy to maintain: each pipeline stage is a separate module that can be run and tested on its own, automated tests cover $\geq$ 80\% of its lines, and static analysis (Ruff, Pylint) reports no errors & M \medskip \\
NFR-4 & The energy and emissions of every pipeline run are measured with CodeCarbon, logged to MLflow and reported in the model card; a full retraining (embed and train) emits $\leq$ 5\,g CO$_2$eq & S \medskip \\
NFR-5 & CPU, memory, request rate, latency and errors are collected every 15\,s and kept for at least 7 days; input drift is tested on each batch of 200 inputs, with an alert when it is significant ($p < 0.05$) & M \medskip \\
NFR-6 & Available $\geq$ 95\% of the time (health check every minute); after a crash the service restarts by itself within 1\,min, and an alert is raised if it is still down after 2\,min; every invalid input gets an error without crashing the service & M \medskip \\
NFR-7 & Uploaded images are only processed in memory, never written to disk or to the logs; files above 10\,MB or not in JPEG/PNG/WEBP are rejected before decoding;  & M \medskip \\
NFR-8 & On the deployment VM, \texttt{POST /predict} answers within 300\,ms for $\geq$ 95\% of requests, and the service is ready $\leq$ 60\,s after start-up & M \medskip \\
NFR-9 & While the service is starting, restarting or cannot load the model, it answers with a 503 error and its readiness check fails; as it keeps no state, a crash only loses the requests in progress, which clients can retry & M \medskip \\
NFR-10 & Besides the label and the score, each response includes the decision threshold and the model version (the Git commit of the release); each deployed version has a model card with its test metrics per generator and its known limitations & M \medskip \\
NFR-11 & An operator can roll the service back to any released model version, or take it offline, in $\leq$ 10\,min & S \medskip \\
NFR-12 & Results are repeatable and consistent: \texttt{dvc repro} at the same commit reproduces every metric within $\pm$0.001; a retrained model is deployed only if it still meets main model requirements on the test split & M \\
\bottomrule
\end{tabularx}

\smallskip
\noindent The time and memory targets keep a wide margin over our measurements on a laptop CPU (Apple M4), the margins also cover a slower server CPU.

\medskip
\medskip

\newpage

\medskip\noindent\textbf{Challenges.}
The project runs on general-purpose laptops and free-tier cloud, so we favour transfer learning on a small model, and the dataset must fit a free-tier DVC remote ($\sim$15\,GB) --- which Defactify ($\sim$7.5\,GB) does. We use the data under \textbf{CC BY 4.0} (non-commercial academic use, MS COCO), assuming inputs resemble the training distribution (everyday photos vs full AI images) rather than face-swaps or adversarial inputs. The main risks are:
\begin{itemize}[nosep,leftmargin=1.4em]

\medskip

  \item \textbf{shortcut learning} --- the model keying on non-semantic cues such as image shape or format instead of ``AI-ness'' (mitigated by normalising size and re-encoding uniformly);

\medskip

  \item \textbf{class imbalance} of roughly 5:1 fake:real (class weights / generator-stratified sampling and balanced metrics);

\medskip


  \item \textbf{split leakage} from captions shared across train/val/test, which we already observed in EDA (a caption-grouped re-split);

\medskip

  \item \textbf{generator drift}, i.e.\ degraded accuracy on unseen or newer generators (leave-one-generator-out evaluation and drift monitoring).
\end{itemize}



\subsubsection{Dataset card}
\textit{Summary of the full card in \texttt{docs/dataset\_card.md}. Source: \href{https://huggingface.co/datasets/Rajarshi-Roy-research/Defactify_Image_Dataset}{Defactify / MS-COCOAI on Hugging Face}, arXiv 2601.00553. Licence: CC BY 4.0, with the real images under MS COCO's terms.}

\paragraph{Summary.} Defactify pairs 16{,}000 real \textbf{MS COCO} photographs with 80{,}000 AI-generated images produced from the \emph{same captions} by five generators (Stable Diffusion 2.1, SDXL, Stable Diffusion 3, DALL-E 3 and MidJourney v6). As the content is controlled, a detector must learn \emph{how} an image was made, not \emph{what} it shows. Each record has a \texttt{Caption}, the \texttt{Image}, \texttt{Label\_A} (0 real, 1 AI) for the main task, real-vs-AI detection, and \texttt{Label\_B} (0 real, 1--5 generator) for the optional generator attribution. We use a versioned subsample of 800 captions (7{,}734 images), split 70/15/15 by caption.

\paragraph{Considerations} (from our exploratory analysis):


\begin{itemize}[nosep,leftmargin=1.4em]

\medskip

  \item \textbf{Class imbalance:} about 5 AI images per real photo.

\medskip

  \item \textbf{Aspect-ratio shortcut:} almost all AI images are square, against $\sim$4\% of the real photos, so every image is centre-cropped, resized and re-encoded.

\medskip

  \item \textbf{Split leakage:} captions repeat across the original splits, so we re-split by caption.

\medskip

  \item \textbf{Scope and bias:} everyday COCO scenes and five 2022--2024 generators only (no face swaps, adversarial inputs or newer generators); the photos may show people and reflect MS COCO's own collection biases.
\end{itemize}

\newpage

\subsubsection{Model card}
\textit{Following Mitchell et al.\ (2019); summary of the full card in \texttt{docs/model\_card.md}.}

\paragraph{Model.} Built by Team Vera with transfer learning: the frozen image encoder of \textbf{CLIP ViT-B/32} (512-d embeddings) with a standardisation + \textbf{logistic-regression} head (C~$=0.01$, balanced class weights against the $\sim$5:1 imbalance). An image is flagged as AI when $P(\text{AI}) \geq 0.57$, a threshold tuned on validation and stored in the model. It was selected among six baselines and a regularisation sweep tracked with MLflow, and trained on the 5{,}352 training images of our subsample.

\paragraph{Intended use.} Deciding whether an image is real or AI-generated, with a confidence score, for content-moderation back-ends, fact-checkers and developers using the API. Out of scope: face-swap or partial-manipulation forensics, adversarial inputs, generators far from the training set, and any high-stakes decision without human review.

\paragraph{Evaluation.} Balanced accuracy, macro-F1, PR-AUC and recall of the minority real class, reported overall and per generator. On validation: balanced accuracy 0.935, macro-F1 0.884, recall (real) 0.947, PR-AUC (real) 0.933, and recall per generator from 0.85 (SD~3) to 0.98 (DALL-E~3). On generators left out of training, balanced accuracy is 0.898 on average and 0.822 at worst (SD~3), so MR-4 is met. The held-out test split is kept for the final evaluation in a later milestone.

\paragraph{Ethical considerations and caveats.} A detector can be evaded, and false positives may wrongly flag genuine images, so outputs are advisory and need human review. Accuracy is expected to drop on newer generators (monitored for drift and retrained when needed), and the model is only valid for COCO-like, modest-resolution images.


\subsubsection{Project coordination and communication}

We coordinate the work with below tools, chosen so that every task, discussion and change leaves a trace that the whole team, and the teaching staff, can follow.


\paragraph{Jira.} We selected because it is a Shared backlog with assignees and statuses, and its issue keys can be referenced from Git. In our flow we have a ticket per each task, assigned to a member and moved through \emph{To Do} $\rightarrow$ \emph{In Progress} $\rightarrow$ \emph{In Review} $\rightarrow$ \emph{Done}. When a ticket is implemented, a comment summarises the work and links its pull requests.

\paragraph{GitHub.} Every change goes through a pull request that references its Jira ticket and is reviewed by a teammate.

\paragraph{Discord.} The team's chat for messages and meetings, including Day-to-day coordination such as: announcing pull requests, asking for reviews, agreeing on decisions and splitting tasks.

\paragraph{DagsHub.} Shared remote for the versioned data and models, and the MLflow tracking server, that each member have access to it with a personal token. \\


The Jira key ties the tools together, it appears in the branch name, in the pull request description and, by convention, in the commit footers (\texttt{Refs: KAN-12}), so every change can be followed from its story to the code and back. Each member's contribution is visible in the assignees of the board and in the authors of the commits and pull requests on GitHub.

\subsubsection{Selection of the cloud provider}

In Milestone~4 the component is deployed as a REST API that only runs inference: a frozen CLIP ViT-B/32 image encoder (about 88M parameters, $\sim$350\,MB) and a logistic-regression head, on CPU. The host therefore needs around 2\,GB of RAM for the model, PyTorch and the API, plus room for the containers and the monitoring stack, It must also stay available, free of charge, for the whole course.

\begin{table}[h]
\centering
\small
\caption{Free tiers of the providers considered (course deployment guide).}
\label{tab:cloud}
\begin{tabular}{@{}lcccc@{}}
\toprule
\textbf{Provider} & \textbf{Compute} & \textbf{RAM} & \textbf{Disk} & \textbf{Free for} \\
\midrule
\textbf{Virtech (FIB)} & unlimited & 4\,GB & 20\,GB & 1 semester \\
Azure Virtual Machines & 750\,h/month & 1--4\,GB & 128\,GB & 12 months \\
AWS EC2 & 750\,h/month & 1\,GB & 30\,GB & 12 months \\
Google Compute Engine & 750\,h/month & 0.6\,GB & 30\,GB & 12 months + 300\,USD \\
Okteto & unlimited & 3\,GB & 5\,GB & unlimited \\
\bottomrule
\end{tabular}
\end{table}
\medskip

We select \textbf{Virtech}, the FIB's private cloud (based on OpenNebula), which provides one virtual machine per team. It offers the most memory among the free options, while the 1\,GB and 0.6\,GB instances of AWS and Google Cloud are too small for the model and PyTorch and Okteto's 5\,GB disk is tight for container images. It also has no usage limits during the semester, needs no credit card or personal cloud account. Its main limitations are that the VM only lasts one semester.


% ----- 2.2 Milestone 2 -----
\subsection{Milestone 2: Model Building -- Reproducibility}
In Milestone~2 we made the project reproducible: the data, the model and its metrics can be rebuilt from the repository and its DVC remote alone. Four practices work together: the \textbf{Cookiecutter Data Science} structure and \textbf{uv} give everyone the same layout and the same locked environment; \textbf{Git with GitHub Flow} versions the code through reviewed pull requests; \textbf{DVC} versions the data and the models as a pipeline whose outputs are stored on DagsHub; and \textbf{MLflow} records every training run. As a result, one Git commit pins the code, the parameters (\texttt{params.yaml}), the pipeline (\texttt{dvc.yaml}), the fingerprints of every data and model file (\texttt{dvc.lock}) and the environment (\texttt{uv.lock}): \texttt{uv sync} and \texttt{dvc pull} recover exactly those versions, and \texttt{dvc repro} rebuilds them. To check it, we rebuilt the model from the raw images on a second machine, and every training metric matched to within $10^{-5}$.

\subsubsection{Project structure}

The repository follows the \textbf{Cookiecutter Data Science} template (v2, generated with \texttt{ccds} for the package \texttt{mlops\_vera}), Table~\ref{tab:structure} summarises the layout.

\begin{table}[h]
\centering
\small
\caption{Repository layout.}
\label{tab:structure}
\begin{tabularx}{\textwidth}{@{}l Y@{}}
\toprule
\textbf{Path} & \textbf{Contents} \\
\midrule
\texttt{data/} & \texttt{raw}, \texttt{interim}, \texttt{processed} and \texttt{external} data, versioned with DVC \\
\texttt{mlops\_vera/} & Python package with one module per pipeline stage: \texttt{dataset.py} (download), \texttt{features.py} (preprocessing), \texttt{split.py} (caption-grouped split) and \texttt{modeling/} (training and inference). \texttt{config.py} holds all paths and loads \texttt{params.yaml}. \\
\texttt{models/} & Trained models, produced by the pipeline and versioned with DVC \\
\texttt{notebooks/} & Exploratory notebooks \\
\texttt{docs/} & Dataset card and model card \\
\texttt{reports/} & This report, the milestone write-ups and metrics (DVC metrics) \\
\texttt{references/} & Background material \\
\texttt{tests/} & Pytest suite \\
Root files & \texttt{pyproject.toml} and \texttt{uv.lock} (environment); \texttt{dvc.yaml}, \texttt{dvc.lock} and \texttt{params.yaml} (pipeline); \texttt{Makefile}, \texttt{README.md} and \texttt{CONTRIBUTING.md} \\
\bottomrule
\end{tabularx}
\end{table}

The environment is managed with \textbf{uv} because it provides fast, reproducible dependency installation and keeps the project environment consistent across team members. The dependencies are declared in \texttt{pyproject.toml}, and the exact versions are locked in \texttt{uv.lock}, and the development tools (Pylint, Pytest, ...) are kept in a separate \texttt{dev} group, so \texttt{uv sync} rebuilds the same environment on any machine. The \texttt{Makefile} wraps the common tasks (\texttt{make data}, \texttt{make test}, \texttt{make lint}).
\medskip

The structure keeps code, data, models and documentation apart, and since every stage reads its paths from \texttt{config.py} and its hyper-parameters from \texttt{params.yaml}, a stage can be changed or re-run without touching the others.

\subsubsection{Code versioning}
\label{sec:code-versioning}

The code is versioned with Git in a public repository of the course organisation, following GitHub Flow \texttt{main} is always deployable, and every change is implemented in its own branch and merged through a reviewed pull request. The main rules which are documented in \texttt{CONTRIBUTING.md}, are:
\medskip

\begin{itemize}[nosep,leftmargin=1.4em]
  \item \textbf{Branches:} one branch per Jira ticket, named \texttt{<type>/<Jira key>-<description>}, e.g.\ \texttt{feat/ KAN-12-mlflow-tracking}.

\medskip

  \item \textbf{Commits:} small, related changes, with messages following the \emph{Conventional Commits}~1.0.0  (\texttt{type (scope):description}), a body that explains what changed and why, and a Jira ticket \texttt{Refs: KAN-n} footer.

\medskip

  \item \textbf{Pull requests:} titled like a commit or with the Jira key and the story summary, and filled in from a template (why, what is changing, how to test, and a checklist covering tests, credentials, data and documentation). A pull request needs the approval of another member before it can be merged, to make sure changes added to main are correct and complete.

\medskip

  \item \textbf{History:} published history is never rewritten (no force pushes), and  conflicts are solved by merging \texttt{main} into the branch first, to keep the main version safe and request the changes to be merged after an approval.
\end{itemize}

\medskip

To minimise mistakes, GitHub enforces these rules instead of relying on discipline. We set the \texttt{main} to be protected and  only accept changes through a pull request with minimum one team member approval, and cannot be force-pushed or deleted. Pull requests are merged with a merge commit, so each member's commits stay in \texttt{main}, and merged branches are deleted automatically to keep the repository clean. Secret scanning with push protection blocks pushes that contain recognised credential formats, credentials are only kept in git-ignored local files (\texttt{.dvc/config.local}, \texttt{.env}), and \texttt{.gitignore} excludes data, caches and these local environments.

\subsubsection{Data versioning}
\guideline{(e.g., DVC)}

\subsubsection{Experiment tracking}

\paragraph{Tooling.} Training runs are tracked with \textbf{MLflow}, using the tracking server that \textbf{DagsHub} hosts for our repository (\url{https://dagshub.com/AdriSegurao/MLOps-Vera.mlflow}). The same DagsHub project already serves as our DVC remote, so code, data and experiments are browsable in one place, for free and shared by the whole team. Credentials live in a git-ignored \texttt{.env} file (template in \texttt{.env.example}); without it, runs fall back to a local SQLite store (\texttt{./mlflow.db}), so the pipeline and its tests never depend on the network.

\paragraph{Integration with DVC.} Model building is split into two DVC stages. \texttt{embed} passes every preprocessed image through a \emph{frozen} pretrained backbone and stores one embedding per image and split. \texttt{train} fits a standardisation + logistic-regression head on those embeddings, tunes the decision threshold on validation (maximising balanced accuracy, as the model card specifies) and logs the run to MLflow. Keeping the costly forward pass in its own stage means that changing only head parameters re-runs \texttt{train} alone, in seconds. An experiment is a variation of \texttt{params.yaml} launched with \texttt{dvc exp run -S <param>=<value>}: DVC records the exact code, parameters, data hashes and metrics of each experiment as a Git reference (\texttt{dvc exp show}, \texttt{dvc exp push}), while MLflow provides the visual comparison and stores the artefacts. The test split is never read by \texttt{train}; it is reserved for the final evaluation, so that choosing between experiments cannot overfit it.

\paragraph{What each run records.}
\begin{itemize}[nosep,leftmargin=1.4em]
  \item \textbf{Parameters:} the \texttt{embed} (backbone, batch size) and \texttt{train} (C, class weighting, iterations, seed) sections of \texttt{params.yaml}.
  \item \textbf{Metrics} on train and validation: the model-card metrics (balanced accuracy, macro-F1, PR-AUC and recall of the minority real class), plus ROC-AUC, the recall of each generator and the tuned threshold.
  \item \textbf{Traceability tags:} the pinned dataset revision, the DVC hash of the embeddings the head was trained on, and the backbone weights; MLflow adds the Git commit. Any run can therefore be tied to the exact code and data that produced it.
  \item \textbf{Artefacts:} the validation confusion matrix and precision--recall curve, and the fitted model with its input signature. The tuned threshold is baked into the logged model (scikit-learn's \texttt{FixedThresholdClassifier}), so its \texttt{predict()} gives the same decisions as the DVC-versioned model.
\end{itemize}
Runs are named \texttt{<backbone>-logreg-<balanced|unweighted>-C<C>} so they can be told apart at a glance in the MLflow UI.

% TODO: add a screenshot of the DagsHub MLflow run comparison, e.g.
% \begin{figure}[h]\centering
%   \includegraphics[width=\textwidth]{images/mlflow-runs.png}
%   \caption{Baseline runs compared in the DagsHub MLflow UI.}
% \end{figure}

\paragraph{Baseline experiments.} We compared three ImageNet- or web-pretrained backbones used as frozen feature extractors: ResNet-18 and ResNet-50 (torchvision) and the image encoder of CLIP ViT-B/32 (OpenAI weights, via open\_clip). Each was trained with and without class weighting against the $\sim$5:1 fake:real imbalance (C~$=1$), giving six runs, cheap enough to run on a laptop CPU with our $\sim$7.7k-image subsample. Table~\ref{tab:baselines} shows the validation results.

\begin{table}[h]
\centering
\small
\caption{Baseline experiments (validation split, 1{,}254 images: 209 real, 1{,}045 AI).}
\label{tab:baselines}
\begin{tabular}{@{}llccccc@{}}
\toprule
\textbf{Backbone} & \textbf{Class weight} & \textbf{Bal.\ acc.} & \textbf{Macro-F1} & \textbf{PR-AUC (real)} & \textbf{Recall (real)} & \textbf{Train bal.\ acc.} \\
\midrule
\textbf{CLIP ViT-B/32} & \textbf{none} & \textbf{0.916} & \textbf{0.845} & \textbf{0.898} & 0.943 & 0.960 \\
CLIP ViT-B/32 & balanced & 0.911 & 0.829 & 0.896 & \textbf{0.952} & 0.949 \\
ResNet-50 & none & 0.821 & 0.717 & 0.582 & 0.871 & 0.960 \\
ResNet-50 & balanced & 0.816 & 0.712 & 0.563 & 0.866 & 0.964 \\
ResNet-18 & none & 0.811 & 0.686 & 0.660 & 0.904 & 0.873 \\
ResNet-18 & balanced & 0.803 & 0.690 & 0.653 & 0.871 & 0.893 \\
\bottomrule
\end{tabular}
\end{table}

\begin{table}[h]
\centering
\small
\caption{Share of each generator's validation images detected as AI (recall per generator).}
\label{tab:per-generator}
\begin{tabular}{@{}llccccc@{}}
\toprule
\textbf{Backbone} & \textbf{Class weight} & \textbf{SD 2.1} & \textbf{SDXL} & \textbf{SD 3} & \textbf{DALL-E 3} & \textbf{MidJourney} \\
\midrule
CLIP ViT-B/32 & none & 0.856 & 0.952 & 0.789 & 0.943 & 0.904 \\
CLIP ViT-B/32 & balanced & 0.809 & 0.952 & 0.775 & 0.933 & 0.880 \\
ResNet-50 & none & 0.708 & 0.809 & 0.713 & 0.871 & 0.751 \\
ResNet-50 & balanced & 0.689 & 0.794 & 0.722 & 0.876 & 0.746 \\
ResNet-18 & none & 0.684 & 0.713 & 0.641 & 0.794 & 0.756 \\
ResNet-18 & balanced & 0.708 & 0.713 & 0.646 & 0.837 & 0.770 \\
\bottomrule
\end{tabular}
\end{table}

\paragraph{Findings and model selection.}
\begin{itemize}[nosep,leftmargin=1.4em]
  \item \textbf{The backbone dominates.} CLIP improves balanced accuracy by about 10 points over both ResNets and raises the PR-AUC of the real class from 0.56--0.66 to 0.90. It is the only backbone that meets MR-1 (balanced accuracy $\geq 0.85$). This is consistent with prior work showing that frozen CLIP features, learnt from web-scale image--text pairs, transfer well to generated-image detection (Ojha et al., CVPR 2023).
  \item \textbf{A bigger ImageNet backbone barely helps.} ResNet-50 gains one point over ResNet-18 but overfits more (0.96 train vs.\ 0.82 validation balanced accuracy); CLIP shows the smallest gap (0.96 vs.\ 0.92).
  \item \textbf{Class weighting hardly changes the scores.} Tuning the threshold on validation already compensates for the imbalance, so weighted and unweighted heads differ by less than one point, which is within noise for 209 real validation images.
  \item \textbf{SD~3 is the hardest generator} for every backbone (0.79 recall with CLIP), while SDXL and DALL-E~3 are the easiest; per-generator results will be studied further with the leave-one-generator-out experiment (MR-4).
\end{itemize}
We keep \textbf{CLIP ViT-B/32} as backbone and \textbf{balanced class weights}: although the scores barely differ here, weighting makes the head independent of the class ratio, so it can be retrained on data with a different imbalance without retuning.

\paragraph{Regularisation sweep.} With C~$=1$ the CLIP head fits the training set almost perfectly (train ROC-AUC 1.000 vs.\ 0.970 on validation). Since the embeddings are cached, a sweep of C only re-runs \texttt{train} (a few seconds per run); Table~\ref{tab:c-sweep} shows that stronger regularisation narrows the gap and improves every validation metric, with C~$=0.01$ best. It also brings the tuned threshold down from 0.99 to 0.57, i.e.\ the balanced head's probabilities are better calibrated.

\begin{table}[h]
\centering
\small
\caption{Regularisation sweep of the CLIP ViT-B/32 head with balanced class weights (validation split).}
\label{tab:c-sweep}
\begin{tabular}{@{}lcccccc@{}}
\toprule
\textbf{C} & \textbf{Bal.\ acc.} & \textbf{Macro-F1} & \textbf{PR-AUC (real)} & \textbf{Recall (real)} & \textbf{ROC-AUC train / val} & \textbf{Threshold} \\
\midrule
1 & 0.911 & 0.829 & 0.896 & \textbf{0.952} & 1.000 / 0.970 & 0.994 \\
0.3 & 0.916 & 0.850 & 0.900 & 0.938 & 1.000 / 0.973 & 0.919 \\
0.1 & 0.919 & 0.871 & 0.910 & 0.919 & 0.999 / 0.976 & 0.707 \\
0.03 & 0.931 & 0.877 & 0.924 & 0.943 & 0.999 / 0.980 & 0.688 \\
\textbf{0.01} & \textbf{0.935} & \textbf{0.884} & \textbf{0.933} & 0.947 & 0.998 / 0.982 & 0.570 \\
\bottomrule
\end{tabular}
\end{table}

The selected model is therefore \textbf{CLIP ViT-B/32 + logistic regression with balanced class weights and C~$=0.01$}, which meets all four validation targets (MR-1--MR-3); it is the configuration committed in \texttt{params.yaml} and versioned with DVC.

% ----- 2.3 Milestone 3 -----
\subsection{Milestone 3: Model Building -- Quality Assurance}
\guideline{Describe how the different software engineering practices and tools have been applied.}

\subsubsection{Energy efficiency awareness}
\guideline{(e.g., CodeCarbon) Reason about the energy efficiency of your ML component.}

\subsubsection{Static code analysis}
\guideline{(e.g., pylint)}

\subsubsection{Model testing}
\guideline{(e.g., pytest)}

\subsubsection{Data testing}
\guideline{(e.g., Great Expectations, Deepchecks)}

% ----- 2.4 Milestone 4 -----
\subsection{Milestone 4: Model deployment -- API}
\guideline{Describe how you have deployed your model and the technologies used.}

\subsubsection{ML-based component/system architecture}
\guideline{Description of the physical and logical architectures of the system, including diagrams. Design patterns used. Selection and justification of design decisions (e.g., cloud provider).}

\subsubsection{API design}
\guideline{(e.g., FastAPI)}

\subsubsection{API testing}
\guideline{(e.g., FastAPI, pytest)}

% ----- 2.5 Milestone 5 -----
\subsection{Milestone 5: Model deployment -- Model packaging}
\guideline{Describe how the different software engineering practices and tools have been applied.}

\subsubsection{Container and Orchestration}
\guideline{(e.g., Docker, Docker compose)}

\subsubsection{CI/CD}
\guideline{(e.g., GitHub Actions)}

% ----- 2.6 Milestone 6 -----
\subsection{Milestone 6: Monitoring}
\guideline{Describe how the different software engineering practices and tools have been applied.}

\subsubsection{Resource monitoring}
\guideline{(e.g., Grafana, Prometheus)}

\subsubsection{Model performance}
\guideline{(e.g., Alibi Detect)}

\subsubsection{Cycle: Feedback loops and retraining}
\guideline{Description of all the stages of the system pipeline, and how they are connected. Explain how the pipeline can be used to automate the retraining process.}

% ===================== 3. SELF-EVALUATION ==========================
\section{Self-evaluation of the project (retrospective)}
\guideline{To do in the final deliverable.}
\guideline{Explain the main challenges, barriers, and opportunities you encountered during the project. Describe what you have learned, what you would do differently, etc.}
\guideline{Explain what concepts you have used from other courses.}

\end{document}
